\documentclass[../main.tex]{subfiles}

\begin{document}

\raggedright\textbf{Abstract}\\
\justifying\noindent
Autonomous drone navigation in forest environments faces significant challenges due to dense obstacles and the expensive cost of traditional sensing systems (LiDAR: \$1,000--10,000+; stereoscopic cameras: \$300--800), limiting scalable adoption for forest health monitoring applications such as bark beetle detection. This thesis investigates whether Proximal Policy Optimization (PPO) reinforcement learning can enable effective tree obstacle avoidance using only low-cost RGB camera input, with quantitative evaluation of sim-to-real transfer efficiency.

\justifying\noindent
We developed a PPO-based navigation policy trained in AirSim forest simulations and deployed it on a Holybro S500 V2 quadcopter equipped with an NVIDIA Jetson Orin Nano companion computer and Arducam IMX477 RGB camera. The system was evaluated across three tree density configurations (sparse: 6\,m spacing, medium: 4\,m, dense: 2\,m) in both simulation and real-world field trials.

\justifying\noindent
Results demonstrated X\% obstacle avoidance success in sparse simulation environments and X\% in real-world deployment, achieving X\% sim-to-real transfer efficiency. Performance degraded predictably with increasing tree density, with collision rates rising from X\% (sparse) to X\% (dense) in real-world trials.

\justifying\noindent
The findings confirm that PPO-based RGB-only navigation enables cost-effective autonomous forest flight under controlled conditions, providing a practical foundation for scalable forest monitoring systems while highlighting the importance of reward design and environmental complexity in sim-to-real deployment.\vspace{5mm}

\raggedright\textbf{Keywords}\\
\noindent
Proximal Policy Optimization, Autonomous Drone Navigation, Forest Environment Simulation, Reinforcement Learning, Sim-to-Real Transfer, Domain Randomization\vspace{5mm}

\end{document}
